\chapter{\label{sec:korr_kapcs}Korrekció a kapcsolófelületen}

\section{Transzverzális metszés, hibrid dinamikai rendszer}
A közös időpontra korrigáló leképezés könyvbeli levezetése: \cite[2.5.1--2.5.2. szakasz]{dib}. A következő számítás az ottani transzverzális eseménygeometriát részletezi.

A korrekciós mátrix levezetéséhez lásd \cite{kong}.
Legyen $\dot{x}=F(x)$ a folytonos dinamika, $\Sigma=\{x:H(x)=0\}$ az eseményfelület, $x^+=R(x^-)$ a reset leképezés. A megoldásoperátor $\Phi(x,t)$. Feltesszük, hogy $F,H,R$ simák és $H_x(x^*)F(x^*)\neq0$ a vizsgált ütközési pontban.

A referenciapálya a $t_1$ időpontban éri el az $x^*$ eseménypontot, majd az állapota $R(x^*)$-ra ugrik. A közeli kezdeti állapotból indított, perturbált pálya általában más időpontban éri el a felületet: legyen ez $t_1+\delta$. A $\delta$ tehát \emph{eseményidő-eltérés}, amely a perturbációtól függ, és pozitív vagy negatív is lehet. Pozitív értéke későbbi, negatív értéke korábbi ütközést jelent.

A stabilitásvizsgálathoz a két pálya eltérését \emph{azonos időpontban} kell összehasonlítanunk. Az $R_x$ önmagában csak az eseményfelületen vett reset leképezés érzékenységét adná meg; nem venné figyelembe, hogy a megzavart pálya más ideig mozog az esemény előtt és után. Ezért bevezetünk egy olyan $Q$ leképezést, amely a perturbált pálya esemény előtti leírását a referencia $t_1$ idejére vonatkozó esemény utáni leírásba viszi át.

Legyen $x_0=x^*+\Delta x$ az esemény előtti sima dinamika szerint $t_1$-ig továbbvezetett perturbált állapot. Ez \emph{segédállapot}: ha a perturbált ütközés korábban történt ($\delta<0$), akkor a tényleges pálya $t_1$-ben már az esemény utáni dinamikát követi, és nem $x_0$-ban van. A sima folytatás éppen azért szükséges, hogy ugyanazon referenciaidőben adhassuk meg a bemenő perturbációt.

A keresett korrekciós leképezés a \ref{fig:korr_idokorrekcio}. ábra jelöléseivel:
\begin{equation}\label{eq:Q_osszetetel}
Q(x_0)=\Phi\bigl(R(\Phi(x_0,\delta)),-\delta\bigr).
\end{equation}
\begin{figure}[htbp]\centering
% Native vector figure. Font sizes are fixed in the final document, no scaling.
\begin{tikzpicture}[x=1cm,y=1cm,>=Latex,line width=.7pt,
 font=\normalfont\fontsize{10.5}{12.6}\selectfont,
 every label/.style={fill=white,inner sep=1.3pt},
 point/.style={circle,fill,inner sep=1.25pt}]
\coordinate (xp) at (5.2,5.95);
\coordinate (xh) at (2.4,5.6);
\coordinate (x0) at (2.3,2.8);
\coordinate (xs) at (6,1.8);
\coordinate (x5) at (3.7,1.225);
\coordinate (rs) at (8.3,2.375);
\coordinate (x3) at (10.2,2.85);
\coordinate (x4) at (9.2,1.2);
\draw[dashed] (1.4,5.475)--(7.4,6.225);
\node[above,rotate=7] at (5,6.55) {Poincaré-metszet};
\draw[red!70!black] (2.5,.925)--(11.1,3.075)
 node[right,black] {$\Sigma$};
\draw[->] (xp)..controls(4.55,4.55)and(4.55,2.6)..(xs);
\draw (xh)..controls(1.55,4.4)and(1.8,3.45)..(x0);
% Matching Bezier derivatives at x0: (2.3,2.8)-(1.8,3.45)=(2.8,2.15)-(2.3,2.8).
\draw[->] (x0)..controls(2.8,2.15)and(3.1,1.25)..(x5);
\draw[->] (rs)..controls(9.05,3.5)and(9,4.3)..(9,4.8);
\draw (x4)..controls(9.8,1.6)and(10.1,2.1)..(x3);
\draw[->] (x3)..controls(10.6,3.6)and(10.75,4.15)..(10.7,4.65);
\draw[densely dashed,->] (x0)..controls(0,-1.5)and(8,-2)..(x4);
\node[below] at (5.8,-1.2) {$Q(x_0)$};
\draw[densely dashed,->] (10.7,2.65)..controls(10.5,1.8)..(9.65,1.1);
\node[right] at (10.55,1.85) {$-\delta$};
\node at (1.45,4.0) {$t_1$};
\node at (4.45,3.65) {$t_1$};
\node[left] at (2.65,1.85) {$\delta$};
\node[point,label=above:$x_p$] at (xp) {};
\node[point,label=above left:$\hat x$] at (xh) {};
\node[point,label=left:$x_0$] at (x0) {};
\node[point,label=below right:$x^*$] at (xs) {};
\node[point,label=below:{$x_5=\Phi(x_0,\delta)$}] at (x5) {};
\node[point,label=above left:{$R(x^*)$}] at (rs) {};
\node[point,label=above right:$x_3$] at (x3) {};
\node[point,label=below right:$x_4$] at (x4) {};
\end{tikzpicture}

\caption{A referencia- és a perturbált esemény időbeli összehangolása.}
\label{fig:korr_idokorrekcio}
\end{figure}

Az \eqref{eq:Q_osszetetel} összetételt belülről kifelé, három lépésben értelmezzük:
\begin{enumerate}
\item \textbf{Eljutunk a perturbált eseményhez.} Az $x_0$ segédállapotból $\delta$ időt lépünk a sima megoldásoperátorral. Az így kapott $x_5=\Phi(x_0,\delta)$ már a kapcsolófelületen van: ez a perturbált pálya tényleges, esemény előtti állapota a $t_1+\delta$ időpontban.
\item \textbf{Alkalmazzuk a reset leképezést.} Az $x_3=R(x_5)$ pont az esemény utáni állapot, ugyanabban a $t_1+\delta$ időpontban. Ütközésnél például ebben a lépésben fordul meg és csökken a normális sebesség.
\item \textbf{Visszatérünk a közös referenciaidőre.} Az esemény utáni $x_3$ állapotot $-\delta$ idővel továbbvezetjük: $x_4=\Phi(x_3,-\delta)=Q(x_0)$. Ezzel ismét a $t_1$ időpontra vonatkozó állapotot kapunk, amelyet már $R(x^*)$-gal kell összehasonlítani.
\end{enumerate}
A harmadik lépésben szereplő negatív idő \emph{matematikai időkorrekció}, nem fizikailag visszafelé lejátszott ütközés. Ha $\delta>0$, az esemény utáni sima pályát visszafelé hosszabbítjuk meg; ha $\delta<0$, előrefelé vezetjük a referenciaidőig. Az itt vizsgált speciális esetben az esemény előtt és után ugyanaz az $F$ mező érvényes, ezért mindkét helyen ugyanazt a $\Phi$ operátort használjuk. Különböző mezők esetén a belső operátor $\Phi_1$, a külső $\Phi_2$, ahogy a \ref{sec:altalanos_korrekcio}. szakaszban látni fogjuk.

A $\delta$ nem független bemenő adat: a közeli eseménypontot a
\[
H\bigl(\Phi(x_0,\delta(x_0))\bigr)=0,\qquad \delta(x^*)=0
\]
feltétel jelöli ki. Mivel ennek $\delta$ szerinti deriváltja a referenciapontban $H_x(x^*)F(x^*)\neq0$, az implicitfüggvény-tétel lokálisan egyértelmű, sima $\delta(x_0)$ függvényt biztosít. Következésképpen $Q$ a választott esemény közelében jól definiált sima leképezés. Nem azt feltételezzük, hogy a teljes hibrid pálya visszafordítható, hanem csak a sima szakaszok lokális folytatását használjuk.

Perturbáció nélkül $x_0=x^*$ és $\delta=0$, így \eqref{eq:Q_osszetetel} szerint $Q(x^*)=R(x^*)$. Kis perturbációnál a releváns esemény utáni eltérés ezért $Q(x^*+\Delta x)-R(x^*)$. Ennek elsőrendű tagját adja majd a $Q_x(x^*)\Delta x$ szorzat; az eseményidő változásának hatása jelenik meg a korrekciós mátrix $R_x$-en felüli tagjában.

Az eseményig történő elmozdulás első rendben:
\begin{align}
x_5&=x^*+\Delta x+\delta F(x^*)+\mathcal O(2),\\
0=H(x_5)&=H_x(x^*)\Delta x+\delta H_x(x^*)F(x^*)+\mathcal O(2),\\
\delta&=-\frac{H_x(x^*)\Delta x}{H_x(x^*)F(x^*)}+\mathcal O(2).
\end{align}
Az $\mathcal O(2)$ jelölés a perturbáció és az eseményidő-eltérés legalább másodrendű tagjait foglalja össze. A reset leképezést és az időkorrekciót is sorba fejtve:
\begin{align}
Q(x_0)&=R(x^*)+R_x(x^*)\Delta x\notag\\
&\quad+\delta\bigl[R_x(x^*)F(x^*)-F(R(x^*))\bigr]+\mathcal O(2).
\label{linearizaltx4}
\end{align}
Ezért
\begin{equation}\label{linearizaltlekepezes}
Q(x_0)=R(x^*)+Q_x(x^*)\Delta x+\mathcal O(\|\Delta x\|^2),
\end{equation}
ahol a korrekciós, más néven saltation-mátrix
\begin{equation}\label{eq:Qxmatrix}
Q_x=R_x+\frac{\bigl[F(R(x^*))-R_xF(x^*)\bigr]H_x}{H_xF(x^*)}.
\end{equation}
Itt $H_x$ egy $1\times n$ sorvektor; a számlálóban oszlop- és sorvektor diadikus szorzata áll, tehát $Q_x\in\mathbb R^{n\times n}$.

\section{Példa: gerjesztett ütköző oszcillátor}
\examplebody

A \eqref{eq:Qxmatrix} képletet most háromdimenziós állapottérben, tényezőnként számítjuk ki. Használjuk az $x=(u,v,s)^{\mathsf T}$ állapotot: $u$ az elmozdulás, $v$ a sebesség, $s$ pedig az időt követő állapotváltozó ($\dot s=1$). A csillapítás $\zeta$, a visszapattanási tényező $r$:
\[
F(x)=\begin{pmatrix}v\\-u-2\zeta v+\cos(\omega s)\\1\end{pmatrix},\qquad
H(x)=u-\sigma,\qquad R(x)=\begin{pmatrix}u\\-rv\\s\end{pmatrix}.
\]
A fal helyzete $\sigma$; ezt megkülönböztetjük a $\delta$ időkorrekciótól. Az ütközés előtti eseménypont legyen $x^*=(\sigma,v^-,s^*)^{\mathsf T}$, az utána következő állapot pedig $R(x^*)=(\sigma,-rv^-,s^*)^{\mathsf T}$. A felső $-$ és $+$ jel az esemény előtti és utáni értéket jelöli, nem hatványt.


A reset leképezés és az eseményfüggvény deriváltját elemenként felírva:
\[
R_x=\begin{pmatrix}1&0&0\\0&-r&0\\0&0&1\end{pmatrix},\qquad
H_x=\begin{pmatrix}1&0&0\end{pmatrix}.
\]
A \eqref{eq:Qxmatrix} képlet minden tagját az eseménypontban értékeljük ki; röviden $F^-=F(x^*)$, $F^+=F(R(x^*))$. Először ezeket a vektorokat számítjuk ki, majd közvetlenül behelyettesítjük a korrekciós képletbe.


Az ütközéskor a sebesség előjelet és nagyságot vált, míg $u$ és $s$ nem ugrik. Emiatt
\[
a^-=-\sigma-2\zeta v^-+\cos(\omega s^*),\qquad
 a^+=-\sigma+2\zeta r v^-+\cos(\omega s^*),
\]
\[
F^-=\begin{pmatrix}v^-\\a^-\\1\end{pmatrix},\qquad
F^+=\begin{pmatrix}-rv^-\\a^+\\1\end{pmatrix}.
\]
Bár ugyanaz az $F$ függvény írja le a szabad mozgást az ütközés előtt és után, különböző állapotokban értékeljük ki, ezért általában $F^+\neq F^-$.

A konkrét tényezőkkel a teljes behelyettesítés:
\[
Q_x=
\begin{pmatrix}1&0&0\\0&-r&0\\0&0&1\end{pmatrix}
+\frac{
\left[
\begin{pmatrix}-rv^-\\a^+\\1\end{pmatrix}
-\begin{pmatrix}1&0&0\\0&-r&0\\0&0&1\end{pmatrix}
\begin{pmatrix}v^-\\a^-\\1\end{pmatrix}
\right]\begin{pmatrix}1&0&0\end{pmatrix}}
{\begin{pmatrix}1&0&0\end{pmatrix}\begin{pmatrix}v^-\\a^-\\1\end{pmatrix}}.
\]
A nevező egy skalár; a számlálóban az oszlopvektor és a sorvektor külső szorzata áll. A következő lépésekben ezeket külön kiszámítjuk.


Először a reset leképezés deriváltját szorozzuk az ütközés előtti vektormezővel:
\[
R_xF^-=\begin{pmatrix}1&0&0\\0&-r&0\\0&0&1\end{pmatrix}
\begin{pmatrix}v^-\\a^-\\1\end{pmatrix}
=\begin{pmatrix}v^-\\-ra^-\\1\end{pmatrix}.
\]
Ezt vonjuk ki az esemény utáni mezőből:
\[
b:=F^+-R_xF^-=\begin{pmatrix}-(1+r)v^-\\a^++ra^-\\0\end{pmatrix}.
\]
A nevező pedig
\[
H_xF^-=\begin{pmatrix}1&0&0\end{pmatrix}
\begin{pmatrix}v^-\\a^-\\1\end{pmatrix}=v^-.
\]
Ez az eseményfelület normális irányú sebessége. A transzverzalitási feltétel itt pontosan $v^-\neq0$.


Az oszlopvektort jobbról szorozzuk a sorvektorral:
\[
bH_x=
\begin{pmatrix}-(1+r)v^-\\a^++ra^-\\0\end{pmatrix}
\begin{pmatrix}1&0&0\end{pmatrix}
=\begin{pmatrix}
-(1+r)v^-&0&0\\a^++ra^-&0&0\\0&0&0
\end{pmatrix}.
\]
Általában $(bH_x)_{ij}=b_i(H_x)_j$. Itt csak az első oszlop nem nulla, mert $H_x=(1,0,0)$: az eseményidő elsőrendű változását közvetlenül az elmozdulás perturbációja határozza meg.
A skalár nevezővel való osztás, majd $R_x$ hozzáadása:
\[
\begin{aligned}
Q_x&=\begin{pmatrix}1&0&0\\0&-r&0\\0&0&1\end{pmatrix}
+\begin{pmatrix}-(1+r)&0&0\\(a^++ra^-)/v^-&0&0\\0&0&0\end{pmatrix}\\[2mm]
&=\begin{pmatrix}-r&0&0\\(a^++ra^-)/v^-&-r&0\\0&0&1\end{pmatrix},\qquad v^-\neq0.
\end{aligned}
\]
A gyorsulásokat behelyettesítve a középső sor első eleme tovább egyszerűsíthető, mert a csillapítási tagok kiesnek:
\[
a^++ra^-=(1+r)\bigl[-\sigma+\cos(\omega s^*)\bigr].
\]


A $Q_x$ az azonos referenciaidőben értelmezett perturbációkat kapcsolja össze. Ha $\Delta x=(\Delta u,\Delta v,\Delta s)^{\mathsf T}$, akkor első rendben
\[
\Delta Q=
\begin{pmatrix}-r&0&0\\(a^++ra^-)/v^-&-r&0\\0&0&1\end{pmatrix}
\begin{pmatrix}\Delta u\\\Delta v\\\Delta s\end{pmatrix}
=\begin{pmatrix}
-r\,\Delta u\\
\dfrac{a^++ra^-}{v^-}\Delta u-r\,\Delta v\\
\Delta s
\end{pmatrix}.
\]
Ez magyarázza, hogy az első diagonális elem $-r$, miközben $R_x$ első diagonális eleme $1$: a $Q_x$ a reset leképezés mellett az eltérő eseményidő korrekcióját is tartalmazza. Az $(a^++ra^-)/v^-$ elem azt írja le, hogyan okoz az elmozdulás perturbációja a módosult ütközési időn keresztül sebességeltérést.

A $v^-\to0$ határesetben a transzverzális levezetés érvényét veszti; ezt a grazing-fejezet tárgyalja. Még ha a számláló is nullává válik, a fenti képlet önmagában nem igazol érvényes korrekciós mátrixot az érintési pontban.
\exampleend

\section{Transzverzális metszés, általános eset}\label{sec:altalanos_korrekcio}

Korrekció a kapcsolósíkon. Az általános esetben:
\[
{\dot{x}} = \begin{cases}
   {\ F}_1(x), & \text{ha } \; H({ x})>0 \\
   {\ F}_2(x), & \text{ha } \; H({ x})<0,
 \end{cases}
\]

\begin{figure}
	\centering
	% All labels have the same final 10.5 pt font as the other chapter-3 drawings.
\begin{tikzpicture}[x=1cm,y=1cm,>=Latex,line width=.7pt,
 font=\normalfont\fontsize{10.5}{12.6}\selectfont,
 every label/.style={fill=white,inner sep=1.3pt},
 point/.style={circle,fill,inner sep=1.25pt}]
\coordinate (xm) at (4,1.3);
\coordinate (xp) at (8.4,2.1);
\draw (2.5,.1)--(7,3.7) node[above] {$\Sigma=\{H=0\}$};
\draw[->] (2,4.5)..controls(1.4,2.7)and(2.4,1.9)..(xm);
\node[left] at (1.25,3.2) {$F_1$};
\draw[dashed,->] (xm)..controls(5.3,.1)and(7,.2)..(xp);
\node[below] at (6.1,.55) {$R$};
\draw[->] (xp)..controls(9.1,2.6)and(9.2,3.8)..(9.0,4.5);
\node[right] at (9.6,3.3) {$F_2$};
\node[point,label=below left:$x^-$] at (xm) {};
\node[point,label=below right:{$x^+=R(x^-)$}] at (xp) {};
\end{tikzpicture}

	\caption{$R(x)$ nem kapcsolósíkra képez le.}\label{fig:501}
\end{figure}

ahol ${x}\rightarrow R({x}),\; \text{ha}\; H(x)=0.$ Továbbá $R(x)$-nek nem kötelező a kapcsolósíkra leképeznie (lásd \ref{fig:501}. ábra).\\

 A $Q(x)$ korrekciós leképezés alatt a következő érthető: (lásd \ref{fig:502}. ábra)
\noindent
 $\Sigma = \lbrace x,\; H(x)=0 \rbrace,\\
 x^*=\Phi_1(x_p,t),\\
 x_0=\Phi_1{(\hat{x},t_1)} \rightarrow x_2=\Phi_1(x_0,\delta),\\
 x_4=\Phi_2(x_3,-\delta),\\
 x_0 \rightarrow x_4=Q(x_0)=\Phi_2(\underbrace{R(\overbrace{\Phi_1(x_0,\delta)}^{x_2})}_{x_3},-\delta)$

\begin{figure}[ht]
\centering
\begin{tikzpicture}[x=1cm,y=1cm,>=Latex,line width=.7pt,
 font=\normalfont\fontsize{10.5}{12.6}\selectfont,
 every label/.style={fill=white,inner sep=1.3pt},
 point/.style={circle,fill,inner sep=1.25pt}]
\coordinate (xp) at (4,5.9);
\coordinate (xh) at (2,5.6);
\coordinate (x0) at (3.15,3.15);
\coordinate (xs) at (4.7,2.56);
\coordinate (x2) at (4,2);
\coordinate (rs) at (7.8,3.1);
\coordinate (x3) at (8.8,1.6);
\coordinate (x4) at (7.1,.8);
\draw (2.5,.8)--(7,4.4) node[above] {$\Sigma$};
\draw[->] (xp)..controls(3.55,4.55)and(3.9,3.2)..(xs);
\draw (xh)..controls(2.05,4.2)and(2.7,3.6)..(x0);
\draw[->] (x0)..controls(3.5,2.8)and(3.65,2.35)..(x2);
\node[left] at (3.6,4.25) {$\Phi_1$};
\node[right] at (3.9,2.8) {$\delta$};
\draw[dashed,->] (xs)..controls(5.5,1.8)and(6.9,2.35)..(rs);
\draw[dashed,->] (x2)..controls(5.4,.5)and(7.6,.65)..(x3);
\draw[->] (rs)..controls(8.7,3.8)and(8.8,4.5)..(8.7,5.0);
\draw (x4)..controls(7.9,.65)and(8.5,.95)..(x3);
\draw[->] (x3)..controls(9.45,2.4)and(9.8,3.2)..(9.8,4.0);
\node[right] at (10.1,3.05) {$\Phi_2$};
\draw[densely dashed,->] (x0)..controls(1,-1.8)and(6,-1.6)..(x4);
\node[below] at (4.5,-.85) {$Q(x_0)$};
\draw[densely dashed,->] (8.6,1.0)..controls(8.05,.25)..(7.35,.35);
\node[below] at (8,.55) {$-\delta$};
\node[point,label=above:$x_p$] at (xp) {};
\node[point,label=above:$\hat x$] at (xh) {};
\node[point,label=left:$x_0$] at (x0) {};
\node[point,label=above:$x^*$] at (xs) {};
\node[point,label=below right:$x_2$] at (x2) {};
\node[point,label=above left:{$R(x^*)$}] at (rs) {};
\node[point,label=below right:{$x_3=R(x_2)$}] at (x3) {};
\node[point,label=below left:$x_4$] at (x4) {};
\end{tikzpicture}

\caption{Korrekció származtatása ha $R(x)$ nem kapcsolósíkra képez le.}\label{fig:502}
\end{figure}

 $Q(x)$ linearizáltjára adódik, hogy:
 \begin{equation}
  Q_x(x^*)=R_x(x^*)+\frac{(F_2(R(x^*))-R_x(x^*)F_1(x^*))H_x(x^*)}{H_x(x^*)F_1(x^*)}\label{eq:general_saltation}
 \end{equation}
  és $H_x(x^*)F_1(x^*)\neq 0$ transzverzális metszés.

\exampleheading{Példa: végállásütközés és motorerő-váltás.}
Egy vízszintes sínen mozgó, $m$ tömegű szánt motor hajt egy merev végállás felé. A végállástól mért távolság $q\geq0$, a sebesség $v=\dot q$; az állapot $x=(q,v)^{\mathsf T}$. A rövid eseménykörnyezetben a súrlódást elhanyagoljuk, a motorerőt az egyes módokban állandónak vesszük. Beérkezés előtt az erő $u_1$, az ütközés érzékelése után a vezérlő $u_2$ erőre vált. A gyorsulások $a_1=u_1/m$ és $a_2=u_2/m$.

Az ütközés előtti és utáni mező, az eseményfüggvény és a reset leképezés:
\[
F_1(q,v)=\begin{pmatrix}v\\a_1\end{pmatrix},\qquad
F_2(q,v)=\begin{pmatrix}v\\a_2\end{pmatrix},\qquad
H(q,v)=q,\qquad R(q,v)=\begin{pmatrix}q\\-rv\end{pmatrix},
\]
ahol $0<r\leq1$ a visszapattanási tényező. A motorerő véges váltása a gyorsulást módosítja; a pillanatszerű sebességugrást a kontaktus impulzusa okozza.

Ebben a hibrid példában az 1 és 2 index az esemény előtti és utáni \emph{üzemmódot} jelöli. A szán ütközés után is a megengedett $q\geq0$ oldalon van: a módot a vezérlő tárolja, nem kizárólag $H$ előjele határozza meg. A \eqref{eq:general_saltation} korrekciós képlet erre is érvényes: az esemény előtti és utáni megfelelő mezőt kell behelyettesíteni. Ez a könyv általános hibrid eseménykezelésének egyszerű alkalmazása \cite[2.5.2. szakasz]{dib}.


Legyen $x^-=(0,v^-)^{\mathsf T}$, ahol $v^-<0$. Ekkor
\[
x^+=\begin{pmatrix}0\\-rv^-\end{pmatrix},\qquad
R_x=\begin{pmatrix}1&0\\0&-r\end{pmatrix},\qquad
H_x=\begin{pmatrix}1&0\end{pmatrix}.
\]
A mezők az esemény megfelelő oldalain:
\[
F_1(x^-)=\begin{pmatrix}v^-\\a_1\end{pmatrix},\qquad
F_2(x^+)=\begin{pmatrix}-rv^-\\a_2\end{pmatrix},\qquad
H_xF_1(x^-)=v^-\neq0.
\]
Az ütközés tehát transzverzális. A számláló vektorkülönbsége
\[
F_2(x^+)-R_xF_1(x^-)
=\begin{pmatrix}-rv^-\\a_2\end{pmatrix}
-\begin{pmatrix}1&0\\0&-r\end{pmatrix}\begin{pmatrix}v^-\\a_1\end{pmatrix}
=\begin{pmatrix}-(1+r)v^-\\a_2+ra_1\end{pmatrix}.
\]


\[
\begin{aligned}
Q_x&=\begin{pmatrix}1&0\\0&-r\end{pmatrix}
+\frac1{v^-}\begin{pmatrix}-(1+r)v^-\\a_2+ra_1\end{pmatrix}
\begin{pmatrix}1&0\end{pmatrix}\\[1mm]
&=\begin{pmatrix}1&0\\0&-r\end{pmatrix}
+\begin{pmatrix}-(1+r)&0\\(a_2+ra_1)/v^-&0\end{pmatrix}
=\boxed{\begin{pmatrix}-r&0\\(a_2+ra_1)/v^-&-r\end{pmatrix}}.
\end{aligned}
\]
A $R_x$ önmagában nem adná meg ezt az eredményt: az eseményidő eltérése miatt a két pálya más ideig mozog a két különböző motorerő mellett.


Válasszuk az $m=1$ kg, $u_1=-1$ N, $u_2=2$ N, $r=0{,}5$, $v^-=-1$ m/s értékeket. Így $a_1=-1$ m/s$^2$, $a_2=2$ m/s$^2$, $v^+=0{,}5$ m/s, és
\[
Q_x=\begin{pmatrix}-0{,}5&0\\-1{,}5\,\mathrm{s}^{-1}&-0{,}5\end{pmatrix}.
\]
Ha a referenciaütközés időpontjában a beérkező, simán továbbvezetett pálya helyzeteltérése $\Delta q=1$ mm, sebességeltérése pedig $\Delta v=0$, akkor az azonos referenciaidőre korrigált esemény utáni eltérés első rendben
\[
\begin{pmatrix}\Delta q^+\\\Delta v^+\end{pmatrix}
=Q_x\begin{pmatrix}0{,}001\,\mathrm m\\0\end{pmatrix}
=\begin{pmatrix}-0{,}0005\,\mathrm m\\-0{,}0015\,\mathrm{m/s}\end{pmatrix}.
\]
A távolabbról érkező szán körülbelül $\delta=-\Delta q/v^-=0{,}001$ s-mal később ütközik. Ezalatt még a beérkező mód szerint gyorsul, majd a visszapattanás után a közös referenciaidőre kell visszavezetni. E két hatást együtt tartalmazza a $-1{,}5\,\mathrm{s}^{-1}$ elem. A negatív $\Delta q^+$ itt a közös időre visszavezetett \emph{segédállapot} eltérése, nem tényleges falba hatolás; az esemény utáni sima mezőt lokálisan visszafelé is használjuk.

A képlet esemény előtti és utáni perturbációkat köt össze, nem a teljes repülési szakasz érzékenységét. Egy teljes mozgás vagy szabályozási ciklus vizsgálatában a folytonos variációs mátrixok közé kell illeszteni. Ha $v^-=0$, a transzverzalitás megszűnik, és ez a levezetés nem alkalmazható.
\exampleend

\subsection*{A képlet fontos speciális esetei}
 \exampleheading{Példa 1.} Filippov-rendszer sliding nélkül, azaz $F_1|_\Sigma \neq F_2|_\Sigma$, valamint $R(x)=x$, azaz önmagába leképezés. \\
 Ilyenkor:
 \[Q_x=I+\frac{(F_2-F_1)H_x}{H_xF_1}, \; \text{mert} \; R_x(x)=I.\]
\exampleend

\exampleheading{Példa 2.} Ütközéses dinamikai rendszer, ahol $F_1=F_2=F$.
Így a derivált:
\[
Q_x(x^*)=R_x(x^*)+\frac{(F(R(x^*))-R_x(x^*)F(x^*))H_x(x^*)}{H_x(x^*)F(x^*)}
\]
\exampleend

\exampleheading{Példa 3.} Filippov-rendszer sliding mozgással, azaz

 \[
 F_s=F_{12}=(1-\alpha)F_1+\alpha F_2, \; \text{és } \; \alpha=\frac{H_x F_1}{H_x(F_1-F_2)}.
 \]
 Ekkor
\[
	Q_x=I+\frac{(F_{12}-F_1)H_x}{H_xF_1}
\]

\begin{figure}[ht]
	\centering
	\begin{tikzpicture}[x=1cm,y=1cm,>=Latex,line width=.7pt,
 font=\normalfont\fontsize{10.5}{12.6}\selectfont,
 every label/.style={fill=white,inner sep=1.3pt},
 point/.style={circle,fill,inner sep=1.25pt}]
\coordinate (xs) at (5,1.7);
\draw (4.3,.6)--(7,4.85) node[above] {$\Sigma$};
\draw[->] (2,4.5)..controls(.9,1.7)and(2.8,.5)..(xs);
\node[left] at (.9,3.25) {$F_1$};
\draw[very thick,->] (xs)--(6.4,3.9);
\node[right] at (5.9,2.9) {$F_s$};
\node[point] at (xs) {};
\node at (6.85,1.05) {sliding};
\end{tikzpicture}

	\caption{Filippov-rendszer sliding mozgással}\label{fig:503}
\end{figure}
\exampleend

